\section{Sūtrasthāna 6: Seasonal living}

% K 006r IMG_0141.JPG" line 1.
% H dscn2986 fol 010.jpg upper line 6
% N A 1267_11_06.jpg upper line 2

\section{Translation}


\begin{translation}    
  
\item [1] 
Now, we will explain the seasonal living (\emph{ṛtucaryā}).
% K has ṛtucarya, which is also possible as a neuter noun (cf. bhikṣācarya, brahmacarya, dauścarya, etc.) But the other witnesses and vulgate have rṭucaryā. 


\item [2] 
Venerable time is indeed % kālo hi bhagavān
% Ḍalhaṇa: hi upadarśane, 
self-existent, % svayaṃbhur
without beginning, middle or end. % anādimadhyanidhano
%
The loss or abundance of the tastes % rasasya madhurādeḥ, rasaśabdenātra rasavaddravyam ucyate'bhedopacārād ity anye
and people's life and death % jīvitamaraṇe ca
depend on it. % atra ... āyatte (agreeing with jīvitamaraṇe)
It is called time (\emph{kāla}) because % iti kālaḥ
it does not disappear % sa ... na līyate
even for a tiny instant (\emph{kalā}). % sūkṣmām api kalān: accusative of extent?: Ḍalhaṇa saḥ kālaḥ, sūkṣmām api stokām api, kalāṃ bhāgaṃ, na līyate gatimattvāt śliṣṭo na bhavati; 
It is called time because % % iti kālaḥ
it brings together or impels living beings.%
    \footnote{According to Ḍalhaṇa, this derivation (\emph{nirukti}) of word time (\emph{kāla}) is based on its action of making all beings into one mass by destroying them (\emph{saṃharaṇād ekarāśīkaroti}) or impels them towards death (\emph{kālayati mṛtyusamīpaṃ nayati}).}

\item [3] 
The venerable sun % bhagavān ādityo 
belongs to time in the form of a year % tasya samvatsarātmano

\end{translation} 
